% day4_seg3_maint/aug_state_vec.tex
\begin{frame}
Augmenting the State Vector
 \begin{itemize}
 \item State Vector Augmentation: Expanding the size of the state vector
\begin{itemize}
	\item i.e., adding new states
\end{itemize}
 \item Why?
\begin{itemize}
	\item new navaids with their own error dynamics (e.g., GPS)
	\item correlated (non-white) process noise
	\item correlated (non-white) measurement noise
\end{itemize}
\end{itemize}
\end{frame}
\begin{frame}
Augmenting the State Vector
\begin{itemize}
  \item Why? (continued)
\begin{itemize}
  \item Correlation $\Rightarrow$ dynamics
\begin{itemize}
	\item No dynamics in white noise
\end{itemize}
  \item If dynamics are present they must be modeled
  \item Otherwise, dynamic effects will be allocated to other states
\end{itemize} 
\end{itemize} 
\end{frame}
\begin{frame}
Augmenting the State Vector
\begin{itemize}
 \item How?
\begin{itemize}
	\item model non-white noise as result of feeding white noise through a linear system
	\item new states are outputs of new integrators
\end{itemize}
  \item State transition matrix $\Phi$ and error covariance matrix $P$ grow as $j^2$, where $j$ is number of new states
  \item Geometry matrix $H$ grows as $m\cdot j$, where $m$ is number of measurements and $j$ is number of new states
\end{itemize} 
\end{frame}
\begin{frame}
Augmenting the State Vector\\
Example: Ship Inertial Navigation Systems (SINS)
\begin{itemize}
  \item Long mission times
\begin{itemize}
	\item Especially submarines
\end{itemize}
  \item Position fixes insufficient to calibrate gyro bias drift w/o further modeling
\end{itemize} 
\end{frame}
\begin{frame}
Augmenting the State Vector\\
Example: Ship Inertial Navigation Systems (SINS)
\begin{itemize}
  \item Gyro bias drift due to slowly changing bias
  \item Can be modeled as random walk
  \item Random walk modeled as integrated white noise
  \item Kalman filter needed to estimate drift 
\end{itemize} 
\end{frame}
\begin{frame}
Augmenting the State Vector\\
Example: Ship Inertial Navigation Systems (SINS)
\begin{itemize}
  \item Assume platform tilts are kept close to zero through velocity measurements damping the Schuler oscillations
  \item North-West-Up (NWU) coordinate frame (or West-South-Up)
\end{itemize} 
\end{frame}
\begin{frame}
Augmenting the State Vector\\
Example: Ship Inertial Navigation Systems (SINS)
\begin{itemize}
  \item Two angular position errors in radians:  $\psi_N$ and $\psi_W$
\begin{itemize}
	\item Note $\psi_N$ is a small angular rotation about the north axis and  thus induces a west position error.
    \item Similarly, $\psi_W$ is a small rotation about the west axis and thus induces a south position error.
\end{itemize}
  \item $\theta_U =\psi_Z - \psi_N \tan(\lambda)$, where $\psi_Z$ is platform azimuth error and $\lambda$ is the latitude.
\end{itemize} 
\end{frame}
\begin{frame}
Augmenting the State Vector\\
Example: Ship Inertial Navigation Systems (SINS)
\begin{itemize}
  \item Error Dynamics:
\begin{align*}
\dot{\psi}_N - \Omega_U\psi_W &= \delta g_N\\
\dot{\psi}_W + \Omega_U \psi_N - \Omega_N\theta_U &= \delta g_W\\
\dot{\theta}_U + \Omega_N \psi_W &= \delta g_U
\end{align*}
where $\delta g_N,\delta g_W,\delta g_U$ are non-white gyro bias drift errors and $\Omega_N=\Omega\cos(\lambda)$ and $\Omega_U=\Omega\sin(\lambda)$ are components of earth rate
\end{itemize} 
\end{frame}
\begin{frame}
Augmenting the State Vector\\
Example: Ship Inertial Navigation Systems (SINS)
\begin{itemize}
  \item Since $\delta g_N,\delta g_W,\delta g_U$ are non-white, the question arises as to what kind of model is appropriate
  \item Empirical data indicates that gyro biases exhibit random walk behavior over time
  \item Since random walk is the integral of white noise, white noise can be viewed (in some sense) as the derivative of random walk
\end{itemize} 
\end{frame}
\begin{frame}
Augmenting the State Vector\\
Example: Ship Inertial Navigation Systems (SINS)
\begin{itemize}
  \item Hence, the following dynamics of random walk apply:
\begin{align*}
\dot{\delta} g_N &= w_N\\
\dot{\delta} g_W &= w_W\\
\dot{\delta} g_U &= w_U
\end{align*}
where $w_N,w_W,w_U$ are white noise processes.
\end{itemize} 
\end{frame}
\begin{frame}
Augmenting the State Vector\\
Example: Ship Inertial Navigation Systems (SINS)
\begin{itemize}
  \item We use the random walk equations to augment the state dynamics equations
  \item Thus we now have a six-state model:
\end{itemize} 
\[
\begin{bmatrix}
\dot{\psi}_N\\
\dot{\psi}_W\\
\dot{\theta}_U\\
\dot{\delta} g_N\\
\dot{\delta} g_W\\
\dot{\delta} g_U
\end{bmatrix}
=
\begin{bmatrix}
0&\Omega_U&0&1&0&0\\
-\Omega_U&0&\Omega_N&0&1&0\\
0&-\Omega_N&0&0&0&1\\
0&0&0&0&0&0\\
0&0&0&0&0&0\\
0&0&0&0&0&0
\end{bmatrix}
\begin{bmatrix}
\psi_N\\
\psi_W\\
\theta_U\\
\delta g_N\\
\delta g_W\\
\delta g_U
\end{bmatrix}
+\begin{bmatrix}
0&0&0&0&0&0\\
0&0&0&0&0&0\\
0&0&0&0&0&0\\
0&0&0&1&0&0\\
0&0&0&0&1&0\\
0&0&0&0&0&1
\end{bmatrix}
\begin{bmatrix}
0\\
0\\
0\\
w_N\\
w_W\\
w_U
\end{bmatrix}
\]
\end{frame}
\begin{frame}
Augmenting the State Vector:  Summary
\begin{itemize}
  \item Straightforward: When new, relevant dynamic relations found, simply add the states to the state vector and add the new dynamics to the $F$ and/or $\Phi$ matrix
  \item Use basic models such as low-order systems driven by white noise to model new states
\end{itemize}
\end{frame}
\begin{frame}
Augmenting the State Vector:  Summary
\begin{itemize}
  \item Some Popular Models
\begin{itemize}
	\item random walk
	\item random constant
	\item random ramp
	\item Gauss-Markov (i.e., exponentially correlated)
	\item combinations of the above!
\end{itemize}
\end{itemize} 
\end{frame}

